\documentclass[11pt]{article}
\usepackage[T1]{fontenc}
\usepackage{lmodern}
\usepackage[margin=1in]{geometry}
\usepackage{amsmath,amssymb,amsthm,mathtools}
\usepackage{booktabs,longtable,array,microtype,xcolor}
\usepackage{enumitem,listings,hyperref,fancyhdr}
\hypersetup{colorlinks=true,linkcolor=blue!45!black,urlcolor=blue!50!black,citecolor=blue!45!black,
pdftitle={A 16-variable Diophantine encoding of ZFC consistency},
pdfauthor={Research contribution prepared for Vladimir Reshetnikov},
pdfsubject={A fixed Diophantine equation, proof, and reproducible arithmetic artifacts}}
\setlist{nosep}
\setlength{\emergencystretch}{2em}
\pagestyle{fancy}\fancyhf{}
\fancyhead[L]{\small A Diophantine encoding of ZFC consistency}
\fancyhead[R]{\small ProveIt research}
\fancyfoot[C]{\thepage}
\setlength{\headheight}{14pt}
\newtheorem{theorem}{Theorem}[section]
\newtheorem{lemma}[theorem]{Lemma}
\newtheorem{proposition}[theorem]{Proposition}
\newtheorem{corollary}[theorem]{Corollary}
\theoremstyle{remark}\newtheorem{remark}[theorem]{Remark}
\newcommand{\Z}{\mathbb Z}\newcommand{\N}{\mathbb N}\newcommand{\pos}{\mathbb Z_{>0}}
\newcommand{\Con}{\operatorname{Con}}\newcommand{\ZFC}{\mathsf{ZFC}}
\newcommand{\pc}{\operatorname{popcount}}\newcommand{\bl}{\operatorname{bitlength}}
\newcommand{\snap}{\texttt{cdd582bec13f908df3711fa9e05d1ddac368a5db}}
\lstset{basicstyle=\ttfamily\scriptsize,breaklines=true,columns=fullflexible,
  keepspaces=true,frame=single,rulecolor=\color{black!20},showstringspaces=false}
\title{A 16-variable Diophantine encoding\\of ZFC consistency}
\author{Research contribution prepared for Vladimir Reshetnikov\\and the ProveIt repository}
\date{October 22, 2026}

\begin{document}
\maketitle
\begin{abstract}
We specify a fixed polynomial $P\in\Z[\mathbf x]$ in sixteen positive
integer unknowns such that ZFC proves
\[
 \bigl(\nexists\mathbf x\in\pos^{16}:P(\mathbf x)=0\bigr)
 \quad\longleftrightarrow\quad\Con(\ZFC).
\]
The polynomial has exact degree $156$ and a complete straight-line
evaluation using $73=41\mathrm M+32\mathrm A$ arithmetic operations.
An equivalent system has seven equations, maximum degree $60$, and a
$65=35\mathrm M+30\mathrm A$ common evaluation circuit. We also give
an equation in $44$ unrestricted integer unknowns and an explicitly
listed quadratic system in $83$ integer unknowns; summing the latter's
$46$ squared residuals gives a single quartic equation.

The construction specializes the repository's native universal
Diophantine representation to one fixed input. Its principal new step
replaces the variable-input Pell decoder with the affine factor
$3W+1-3W_0$: an integer unit congruent to $1$ modulo $3$ must be $1$.
The fixed coefficients are completely specified by finite formulas and
standalone code starting from the literal K215 table. Their enormous
expanded numerals are not materialized. The proof composes the inherited
K215 calibration and native compiler theorem with new normalization,
input-elimination, and algebraic arguments. The new composed theorem is
a mathematical proof, not a newly kernel-checked Lean or Rocq endpoint.
No global minimality claim is made.
\end{abstract}

\clearpage
\tableofcontents
\clearpage

\section{Result, scope, and the meaning of simplicity}
\label{sec:result}

Fix the ordinary syntactic consistency sentence
\[
 \Con(\ZFC)\equiv
 \neg\exists n\in\N\;\operatorname{Proof}_{\ZFC}(n,\bot),
\]
for the finite-derivation coding used by the K215 calibration. Changing
to a standard effectively intertranslatable coding does not change the
result: the translations and their correctness can be proved in ZFC.
We encode a refutation of ZFC, not a proof of its consistency and not
a semantic truth predicate.

The source snapshot used throughout is \snap. Its K215 endpoint says
that the exact $215$-state binary machine halts on a blank tape if and
only if ZFC has a finite refutation \cite{K215,Endpoint}. The repository
also contains a complete native Diophantine compiler and several
carefully distinguished operation/degree/witness tradeoffs
\cite{Raw75,Compiler75,Bezout85,Scaled84}. Those are substantial inherited
results; the present work does not rebrand them as new discoveries.

\begin{theorem}[Main construction]\label{thm:main}
Let the fixed integers in Section~\ref{sec:constants} be constructed from
the literal K215 table by the prescribed finite recipe. Let $P$ be the
polynomial in Section~\ref{sec:polynomial}. Then
\[
 \ZFC\vdash
 \left[\exists\mathbf x\in\pos^{16}\ P(\mathbf x)=0\right]
 \longleftrightarrow\neg\Con(\ZFC).
\]
It has exact degree $156$, and the supplied complete evaluator has
$41$ multiplications and $32$ additions or subtractions.
\end{theorem}

The following alternatives are included in the archive. All operation
counts charge a binary addition, subtraction, or multiplication as one
operation, including multiplication by a fixed numeral. Fixed numerals
and reuse of earlier results are free in this arithmetic-circuit model.
The bit lengths and the cost of constructing the coefficients are not
included in these counts.

\begin{center}
\small
\begin{tabular}{@{}lrrrr@{}}
\toprule
Representation & Unknowns & Equations & Degree & Arithmetic\\
\midrule
Scaled product $P$ & $16$ positive & $1$ & $156$ & $73$\\
Unscaled product $U$ & $16$ positive & $1$ & $144$ & $74$\\
Comparison system & $16$ positive & $7$ & $60$ & $65$\\
Direct sum of squares & $16$ positive & $1$ & $120$ & $85$\\
Integer substitution & $44$ integers & $1$ & $312$ & $171$\\
Quadratic system & $83$ integers & $46$ & $2$ & not minimized\\
Sum of quadratic squares & $83$ integers & $1$ & $4$ & not minimized\\
\bottomrule
\end{tabular}
\end{center}

``Simplest'' has no invariant meaning without an objective. Number of
unknowns, degree, circuit length, coefficient height, and proof complexity
give different orderings. Nine-natural-variable universal constructions
exist, with extremely large degrees; this is not a variable-count record
\cite{Jones82,BayerDavid}. Our main choice optimizes a short evaluator
and a readable proof within the supplied repository framework. The
seven-equation version is preferable when a system is acceptable and
the number of arithmetic operations matters more than having one equation.
The quadratic version provides a small-degree system. A concrete
unsatisfiable linear integer system has an effective finite arithmetic
proof of unsatisfiability, so such a system cannot furnish an independent
nonexistence statement for consistent ZFC. We make no global degree,
variable-count, or equation-count minimality claim.

The old $84$-operation construction carries a varying ordinary input.
The present coefficients encode one fixed computation, including its
fixed input. The two operation counts therefore describe different
interfaces. The improvement is a specialization, not an assertion that
the old varying-input universal frontier has been replaced.

\section{The polynomial itself}\label{sec:polynomial}

The sixteen unknowns, all strictly positive for the principal theorem, are
\[
 J,F,\alpha,t,f,h,i,T,s,w,\tau,\eta,\zeta,y,Z,\gamma.
\]
Use the six fixed integer coefficient ports
\[
 \beta=B-1,\qquad K_0,\qquad e=2d,\qquad M_C,\qquad
 \widehat M_F=M_F+B-1,\qquad c_0=1-3W_0.
\]
Here $B=2^d$, $W_0=2^{b+2d}$, and $M_F$ is the \emph{native}
verification mask. The port $\widehat M_F$ is its shifted value.
These are fixed coefficients, not existential variables. In particular
$c_0$ is negative; this does not conflict with positivity of the unknowns.

Define the following polynomial abbreviations:
\begin{align}
 q&=\beta J+1,& X&=wq,&Y&=sq^3,&E&=XY,\label{eq:base}\\
 k&=\eta+\zeta,&c&=kY+\eta,&a&=E+Y,&H&=4a+3,\label{eq:pellbase}\\
 \Delta&=a^2+H,&D&=X+ac+\gamma H,\label{eq:disc}\\
 C&=q-F-Z-\alpha-e,&W&=C-Z,\label{eq:outer}\\
 R&=(q^2-Z-qF)(q^2-1)+(M_C+q\widehat M_F)J,\label{eq:packed}\\
 L&=E(kY),&S&=\Delta i^2c^4,&A_0&=\Delta^2 i^2c^4,\label{eq:auxbase}\\
 V&=c(Tf-1)-Rf^2.\label{eq:V}
\end{align}
Equations~\eqref{eq:base}--\eqref{eq:V} are definitions, not additional
constraints or additional unknowns. In particular $R,C,W$ are computed
polynomials; they are not assumed positive off the zero set.

The seven factors are
\begin{align}
 N_1&=\tau^2-L(L+k), &N_2&=D^2-\Delta c^2,\label{eq:norm12}\\
 N_0&=3W+c_0, &N_3&=A_0(V^2-y^2)+y^2,\label{eq:norm03}\\
 N_4&=k-hE-R, &N_5&=(K_0+w)C+q-F-t(q-1),\label{eq:norm45}\\
 N_6&=f^2-S.\label{eq:norm6}
\end{align}
Our two product polynomials are
\begin{equation}\label{eq:main}
 \boxed{\quad
 U=N_1N_2N_0N_3N_4N_5N_6-1,
 \qquad P=\Delta U.
 \quad}
\end{equation}
The required equation is $P=0$ on $\pos^{16}$. Everything in
\eqref{eq:main} is an ordinary polynomial with integer coefficients.
There is no exponent whose value depends on an unknown. The exponents
occurring in the coefficient recipe, such as $b+2d$, are fixed integers.

The equivalent seven-equation system is especially easy to state:
\begin{equation}\label{eq:system}
 \boxed{\quad W=W_0,\qquad N_1=N_2=N_3=N_4=N_5=N_6=1.\quad}
\end{equation}
The nontrivial assertion is that a product zero really forces these
seven equalities; multiplying factors is not automatically a valid
replacement for a system. Section~\ref{sec:proof} proves this assertion.

\section{From the literal K215 table to a fixed input}\label{sec:machine}

The canonical dense K215 file has SHA-256
\begin{center}\small\ttfamily
fdbcef0825886035c3091c8537a69535bc96f9e265f0ae45863543d40ae70090.
\end{center}
Its working states are $0,\ldots,214$, its separate halt is $215$, and
its alphabet is $\{0,1\}$. The exact file is included, together with
the generator that verifies this hash before using it.

The native tableau compiler expects a first step that writes a flag
and stays in a different nonhalting state. Add four initialization
states $0,1,2,3$; represent the old state $q$ by $q+4$, and use a third
tape symbol $2$. The halt becomes $219$. The six exceptional
initialization transitions are:
\begin{center}
\begin{tabular}{@{}rrrrl@{}}
\toprule
State & Read & Next & Write & Move\\\midrule
0&1&1&2&stay\\
1&2&2&2&left\\
2&1&2&0&left\\
2&0&3&0&right\\
3&0&3&0&right\\
3&2&4&0&stay\\\bottomrule
\end{tabular}
\end{center}
Every old transition is copied under $q\mapsto q+4$. All other entries
are identity stay transitions, including the absorbing halt entries.
This defines all $660$ rows of a $220$-state, three-symbol table.
There are $219$ working states; this normalization is not a smaller
binary Busy Beaver candidate.

\begin{lemma}[Blank initialization]\label{lem:normalize}
For every $u\ge1$, on a tape containing ones exactly at
$-u,\ldots,0$, the normalized machine reaches the blank tape, head
position $0$, and state $4$ after exactly $2u+4$ transitions. Its
subsequent execution is the literal K215 blank execution.
\end{lemma}
\begin{proof}
The first two steps mark the origin and move left. State $2$ erases
the $u$ left-hand ones, then moves right from the next blank. State
$3$ crosses the erased cells, erases the origin flag, and enters state
$4$. Counting these steps gives $2u+4$. The remaining claim is induction
on the copied transition table.
\end{proof}

At ordinary compiler input $x=1$, the marked displacement is $2x=2$,
the initial phase word has three headless input ones, and the head
cell contributes one more one. Thus $u=3$, and initialization takes
ten transitions. This normalized machine accepts every admissible
positive input if K215 halts, and no such input otherwise. No unknown
halting time occurs in its finite table.

\subsection{A smaller normalization by allowing both blank margins}

A folded one-sided simulation is unnecessary. The actual local
tableau rules already permit an initial row
\[
 L^{a_\ell} I^u Q R^{a_r}
\]
with arbitrary positive blank margins. The inherited completeness
presentation fixed $a_r=2$ after choosing a one-sided machine; the
local rules contain no such length restriction.

Given any finite accepting run of the two-sided normalized machine,
let $\ell,r$ be its extreme head coordinates, including initialization.
Choose
\[
 a_\ell\ge\max(2,1-u-\ell),\qquad
 a_r\ge\max(2,r+1).
\]
The width $h_*=a_\ell+u+a_r+2$ can then be any sufficiently large
power of five, by enlarging the left margin. Choose independently
a sufficiently large power-of-five time height $H_*\ge T_*+2$ and
pad with the absorbing halt. This retains all width, height, and
dummy-bit padding required by the inherited compiler.

Place the distinguished Start at the last $I$, in column
$x_S=a_\ell+u$, and End at the first $I$. Their separation is $u-1=2$.
The usual helical extension
\[
 \mathcal T(z,t)=G[z\bmod h_*,\ (t+\lfloor z/h_*\rfloor)\bmod H_*]
\]
has periods $(h_*,-1)$ and $(0,H_*)$. Its cyclic index
$-(z-x_S)-h_*(t-1)$ modulo $h_*H_*$ puts Start at $0$ and End at $2$.
The vertical boundary carries no payload, so the seam preserves the
same local rules. Conversely those rules recover a finite strip,
a unique initialized head, a genuine run without boundary escape,
and a halted head before the next horizontal boundary. This
soundness argument never needed a one-sided tape.

This is a new elementary extension of the existing completeness
argument. Its proof and a two-sided accepting toy tableau check are
included separately; it is not attributed to an existing formal
kernel theorem.

\section{Closed definitions of the coefficients}\label{sec:constants}

The coefficient recipe is fully effective and finite. Its input is
the complete normalized table just defined, not an unspecified
universal-machine index. The supplementary source
\texttt{src/k215\_constants.py} fixes all ordering conventions and
implements the formulas below. A reader can therefore recover exactly
one polynomial, rather than an existentially chosen polynomial.

\subsection{The ordered local window list}

The tile alphabet has vertical and horizontal boundary symbols;
an unphased interior symbol for each tape symbol paired with either
no head or one of the $220$ states; and four initialization-phase
symbols $L,I,Q,R$. In this order it has
\[
 a_* =2+3(220+1)+4=669
\]
members. Enumerate allowed $3\times3$ windows with the prescribed
Start window first, End second, then all other accepted tuples in
lexicographic order. The complete predicate is the literal
\texttt{local\_valid} function. It enforces propagated boundaries,
the phase graph $L\to L,I$, $I\to I,Q$, $Q\to R$, $R\to R$,
the exact machine update, no boundary crossing, and a halted head
just below the next horizontal boundary. Neighbor overlap enforces
the same conditions at every window center.

With the supplied tile order the two distinguished tuples are
\begin{align*}
 \mathrm{Start}&=(1,1,1,666,666,667,223,223,446),\\
 \mathrm{End}&=(1,1,1,665,666,666,2,223,223).
\end{align*}
Let $\mathcal W=(w_0,\ldots,w_{k_*-1})$ be this ordered finite
list. An intentionally simple enumeration inspects
\[
 669^9=26843247761153995367677629
\]
candidate tuples. It is not a feasible routine to run for this machine.
It nevertheless defines a particular finite list and contains no
unbounded search for a computation or a refutation.

The four corner entries do not enter the local predicate. Counting
the five-cell crosses instead gives an exact, quickly evaluable
cardinality without constructing the window list:
\[
 k_*=(896017122+235501665872+11505)669^4
     =47353022820782362885779.
\]
The three summands count boundary, unphased, and phased crosses.
The supplied counting function derives them from the actual transition
table; its derivation and an independently enumerated small fixture
are included in the coefficient note.

\subsection{A finite arithmetic recipe}

Here all exponents and summation indices are fixed data. Write
$A_*=2^{\max(2,\bl(k_*+1))}$. For $r,c\in\{0,1,2\}$ and
$0\le v<a_*$, put
\[
 p(r,c,v)=k_*+3a_*+(2-r)3a_*+(2-c)a_*+v.
\]
Initialize selector coefficients $c_j=1$ for $j<k_*$ and
$\mu=A_*-2$. Enumerate the $9a_*$ triples $(r,c,v)$ lexicographically,
with index $j=1,\ldots,9a_*$. Set $c_{p(r,c,v)}=A_*^j$, add
$A_*^j$ to $\mu$, and add it to each selector coefficient whose
window has value $v$ at position $(r,c)$. Reserve four more powers
$A_*^{9a_*+1},\ldots,A_*^{9a_*+4}$ as anchor coefficients, adding
each to $\mu$ and to every selector coefficient.

Set $m=2\pc(\mu)+12a_*+2$. Until
$m+1\ge k_*+9a_*+5$, add the next unused power of $A_*$ to
$\mu$ and increase $m$ by two. Add
$m+1-k_*-9a_*-4$ consecutive dummy positions starting at
$k_*+12a_*$. Let $\mathcal E_0$ contain the selectors, copy
positions, and dummies. Set $M_* =\max\mathcal E_0+3a_*+1$;
add anchors $M_*,3M_*,9M_*,27M_*$, with the four reserved
coefficients, to obtain $\mathcal E$. Undefined coefficients are zero.
If $E_* =\max\mathcal E$, define
\begin{align*}
 H_c&=E_*+24M_*+3a_*+1,& T_1&=H_c+2E_*+a_*+1,\\
 T_2&=T_1+2E_*+1,&g&=T_2+E_*+1.
\end{align*}
Let $L_*$ be the least power of five at least $g+E_*+1$.
In a formal radix $z$, form
\begin{align*}
 D_C(z)&=z^{3a_*}+z^{H_c+a_*}+z^{8M_*}+z^{24M_*}
       +\sum_e c_e(z^{T_1-e}+z^{T_2-e}),\\
 D_R(z)&=z^{H_c},\\
 M_{F,0}(z)&=\mu(z^{T_1}+z^{T_2})
  +\sum_{r<2,c<3,v<a_*}z^{p(r,c,v)}\\
 &\hspace{18mm}+\sum_{r<3,c<2,v<a_*}z^{H_c+p(r,c,v)}
  +z^{9M_*}+z^{27M_*}.
\end{align*}
Add $z^g$ to $D_C$ precisely when
$2D_C(2)-2^{H_c}\equiv0\pmod5$. After this correction the
difference $2D_C(2)-2^{H_c}$ is nonzero modulo five. Put
\[
 \mathcal M=(|\mathcal E|+2)\left(2\sum_e c_e+6\right).
\]
Choose $b$ to be the least power of five satisfying
$2^b\ge\max(4\mathcal M+8,2\mu+4,16)$, and set
$z_*=2^b$, $d=bL_*$, $B=2^d$. If $e_*$ is the first dummy
position, the final constants are
\begin{align}
 M_C&=B-1-\sum_{e\in\mathcal E,\ e\ne1}z_*^e-2z_*^{e_*},\label{eq:MC}\\
 M_F&=M_{F,0}(z_*)+4, &\widehat M_F&=M_F+B-1,\label{eq:MF}\\
 K_0&=D_C(z_*)+B D_R(z_*),&W_0&=2^{b+2d},\label{eq:K0}\\
 \beta&=B-1,&e&=2d,&c_0&=1-3W_0.\label{eq:c0}
\end{align}
The implementation resolves all finite sums in exactly this order.
The inherited compiler proves, among its other typing and
synchronization facts,
\begin{gather*}
 0<M_C,M_F<B-1,\qquad \pc(M_C)+\pc(M_F)=d,\\
 M_C\equiv2\pmod4,\qquad M_F\equiv4\pmod8.
\end{gather*}
The shift in \eqref{eq:MF} is essential.

\subsection{What is explicit, and what was actually evaluated}

The scalar parts of this recipe have been evaluated exactly:
\begin{center}
\begin{tabular}{@{}lr@{}}
\toprule
Quantity & Exact value\\\midrule
Coefficient radix $A_*$ & $2^{76}$\\
Zero-clause padding steps & $23676511410391181435787$\\
Number of dummy positions & $1$\\
Anchor unit $M_*$ & $47353022820782362895815$\\
Cell width $L_*$ & $5^{36}$\\
Radix-bit parameter $b$ & $5^{35}$\\
Bit width $d=bL_*$ & $5^{71}$\\\bottomrule
\end{tabular}
\end{center}
In particular the following large fixed coefficients are already
available in short closed form:
\begin{equation}\label{eq:exactradix}
 \boxed{\quad B=2^{5^{71}},\qquad
 W_0=2^{\,5^{35}+2\cdot5^{71}}.\quad}
\end{equation}
For the determination of $b$, the target $2\mu+4$ dominates the
competing mass target: $\bl(\mu)=1799414867189729789577713$,
whereas the mass target has bit length at most $458130$.
The exact digit layout then makes $5^{35}$ the least admissible
power of five. These statements use finite counting and integer
arithmetic, with no window enumeration or machine-halting assumption.

The fixed coefficients are specified as exact finite integers by
\eqref{eq:MC}--\eqref{eq:c0} and the completely specified window list.
The archive contains their executable definition, the literal
normalized table, its digest, exact scalar parameters, and the actual marker tuples. It does
not contain expanded decimal or hexadecimal representations of the
K215 coefficients. The default generator deliberately does not start
the astronomical enumeration. Its opt-in generation modes execute
the whole finite recipe in principle.

Thus this is a closed coefficient specification, rather than a
materialized coefficient list. It is substantially more explicit
than ``take the index of a machine searching for a contradiction'';
every table entry, finite predicate, ordering rule, and arithmetic
formula is given. It is also a genuine limitation if one wants a
small human-readable coefficient height. We do not claim such a
height bound or optimality in that metric.

\section{Correctness of the fixed-input reduction}\label{sec:proof}

We prove a reversible reduction to the exact normalized
$85$-operation parent of \cite{Bezout85}. That parent has eighteen
positive witnesses, with a variable ordinary input $x$. In its
notation our $\gamma$ is $\rho+\sigma$. Its variable-input block is
a Pell norm with additional positive witness $\delta$.

\subsection{An affine input guard}

\begin{lemma}\label{lem:unit}
If $U=0$ on positive supplied coordinates, then $W=W_0$ and
$N_1=N_2=N_3=N_6=1$. Moreover $N_4,N_5\in\{1,-1\}$ and
$N_4N_5=1$.
\end{lemma}
\begin{proof}
Every factor of the integer product $1$ is a unit. Since
$N_0=3(W-W_0)+1\equiv1\pmod3$, it cannot be $-1$; hence
$N_0=1$ and $W=W_0$.

For the first norm put $v=XY^2>1$. Then
$L(L+k)=v(v+1)k^2$. The negative Pell equation
$x^2-v(v+1)z^2=-1$ has no positive solution when $v>1$.
Indeed, take a solution with least $z>0$. Its inequalities
$vz<x<(v+1/2)z$ show that multiplication by the inverse unit
$(2v+1)-2\sqrt{v(v+1)}$ produces a solution with
$0<z'=(2v+1)z-2x<z$, a contradiction (take the absolute value
of the new first coordinate if necessary). Thus $N_1=1$.

The integer $\Delta$ is $0$ or $3$ modulo four. Therefore neither
$D^2-\Delta c^2$ nor $f^2-\Delta i^2c^4$ can be $-1$ modulo
four. Finally $A_0$ is a square, so $N_3$ is congruent either
to $y^2$ or to $V^2$ modulo four and cannot be $-1$. The claimed
unit signs follow.
\end{proof}

The affine guard saves one gate over the tempting positive guard
$1+(W-W_0)^2$. The constant $1-3W_0$ is evaluated once as a
fixed coefficient. Omitting the multiplier three would be invalid:
the second unit sign would then allow $W=W_0-2$.

\subsection{Outer bounds before decoding anything}

Because $W=W_0>0$ and $Z>0$, equation~\eqref{eq:outer} gives
$C=Z+W_0>0$ and $F+Z<q$. Also $q\ge B\ge16$. The actual native
mask bounds, with the shifted source mask, imply
\begin{equation}\label{eq:Rbounds}
 (2q-1)(q^2-1)<R<q^4-q^3,
 \qquad 3q+1<R-2<R+2<q^4.
\end{equation}
For example $F+Z\le q-1$ implies
$q^2-Z-qF\ge q+(q-1)Z\ge2q-1$. The upper bound follows
from $F,Z\ge1$ and the mask's bound below $2q^2-2q$.
Consequently
\begin{equation}\label{eq:Ebounds}
 E=wsq^4\ge q^4>R+2,\qquad a=E+Y>R+2.
\end{equation}
These estimates establish $R>0$ before any Pell-index or typed-word
argument. Eliminating a positive packed-index variable without
recovering such a bound would be a genuine soundness error.

\subsection{Restoring the auxiliary quotient coordinates}

From $N_6=1$ we have $f^2=1+\Delta i^2c^4$ and
$f>c^2>2c$. As $R<a<\Delta$,
\begin{equation}\label{eq:gapK}
 A_0-(Rf^2+c)=(\Delta-R)f^2-\Delta-c
 \ge f^2-\Delta-c>0.
\end{equation}
The auxiliary norm reads $A_0V^2-(A_0-1)y^2=1$.
It excludes $V=0$. If $v=|V|>1$, then
$v^2-1=(A_0-1)(y^2-v^2)>0$, so $y\ge v+1$ and
$v\ge2A_0-1$. On the other hand the literal formula for $V$
and $T>0$ give $V>-Rf^2-c>-A_0$. Thus a negative $V$ with
$|V|>1$ is impossible. The values $V=\pm1$ are excluded by
$V\equiv-c\pmod f$ and $f>2c>2$. Hence $V>0$.

Define integers
\begin{equation}\label{eq:restoreoj}
 o=cT-Rf,\qquad j=Tf-R\Delta i^2c^3-1.
\end{equation}
Then $of=V+c>0$ and $jc=V+R>0$, so $o,j>0$ and
\[
 V=of-c=jc-R.
\]
This restores the old auxiliary coordinates with their required
positivity, before invoking the parent universal theorem.

\subsection{The main Pell rank and the remaining signs}

Write $A=a+2$ and define $\chi_A(n),\psi_A(n)$ by
\[
 (A+\sqrt{A^2-1})^n=\chi_A(n)+\psi_A(n)\sqrt{A^2-1}.
\]
Let $P_*=2XY^2+1$. The positive first and main norms classify as
\[
 \tau=\chi_{P_*}(n),\quad k=2\psi_{P_*}(n),\qquad
 D=\chi_A(p),\quad c=\psi_A(p),
\]
with $n,p\ge1$. Put $\epsilon=N_4$. Since $P_*\equiv1\pmod E$,
$2n\equiv R+\epsilon\pmod E$. Equations~\eqref{eq:Rbounds}
and \eqref{eq:Ebounds} give $n\ge(R-1)/2\ge24$.
As $P_*>A$ and $c>k$, Pell monotonicity gives $p>n$; consequently
$c>A\Delta^2$, $c>2p$, and $c>2R$.

At this point we use two precisely identified local lemmas from the
inherited proof, not its full input-decoding theorem. The strong-rank
lemma of \cite{Relaxed} applies with $i_{\rm old}=\Delta i$, because
\[
 (i_{\rm old}c^2)^2=\Delta(f^2-1).
\]
It gives $f=\chi_A(m)$, $c\mid m$, and $m\ge c>2p$.
The signed auxiliary step-down argument in \cite[Section 4]{Reverse89}
then applies to $V=of-c=jc-R>0$, using
$0<R<c/2$ and $0<p<c/2$, and yields
\begin{equation}\label{eq:pR}p=R.\end{equation}
Neither local lemma uses the removed input norm. Their hypotheses
and this application were independently reviewed for this report.

Write $2n=R+\epsilon+vE$ with $v\ge0$. If $v\ge1$, then
$2n>2R$, contradicting $n<p=R$. Thus $2n=R+\epsilon$.
If $\epsilon=-1$, then $R=2n+1$. The duplication identity and
$\chi_A(2)>P_*$ imply
\[
 \psi_A(2n)=2A\psi_{\chi_A(2)}(n)\ge A k.
\]
Since $A>Y+1$, we would have $c=\psi_A(2n+1)>k(Y+1)$,
contrary to $c=kY+\eta$ with $0<\eta<k$.
Therefore $N_4=1$, and Lemma~\ref{lem:unit} gives $N_5=1$.
All seven equalities in \eqref{eq:system} have now been recovered.

\subsection{Recovering a power without using the omitted input block}

This step prevents a circular appeal to the old input theorem.
Write $R=2r+1$, so $n=r+1$ and $r\ge1$.
Elementary Pell growth gives
\[
 \frac ck\ge\frac12
 \left(\frac{(2A-1)^2}{2P_*}\right)^r
 >\frac12\left(\frac{(X+1)^2}{X}\right)^r.
\]
The strict inequality follows by expansion using $a=Y(X+1)$.
Together with $c/k<Y+1$ and $X\ge16$, it implies
$Y>X^r/3$ and
\[
 a=Y(X+1)>X^{r+1}/3>2\cdot4^r=2^R.
\]
Set $E_A(j)=\chi_A(j)-a\psi_A(j)$. It starts with $1,2$,
obeys the Pell recurrence, and satisfies
\begin{equation}\label{eq:powercong}
 E_A(j)\equiv2^j\pmod H.
\end{equation}
Indeed $H=4A-5$, so $2A\,2^{j+1}-2^j\equiv2^{j+2}\pmod H$.
The recurrence therefore gives the same values modulo $H$.
The retained definition of $D$ says $E_A(R)=X+\gamma H$.
Both $X$ and $2^R$ lie strictly between $0$ and $a<H$, so
\begin{equation}\label{eq:Xpower}X=2^R.\end{equation}

Since $q\mid X$, write $q=2^v$. The divisibility
$2^d-1\mid2^v-1$ implies $d\mid v$, so $q=B^N$ for $N\ge1$.
Put $u_0=b+2d$. The fixed marker gives
$2^{u_0}=W_0<q=2^v$; and \eqref{eq:Rbounds} gives $v<R$.
Thus $u_0<R$. Both are odd, so $R\ge u_0+2$.

\subsection{Restoring the three deleted input witnesses}

Define
\begin{align}
 \kappa&=\psi_A(u_0),&\mu_0&=\chi_A(u_0),\notag\\
 \delta&=\frac{\kappa-u_0}{\Delta},&
 \rho&=\frac{\mu_0-a\kappa-W_0}{H},&
 \sigma&=\gamma-\rho.\label{eq:restoreinput}
\end{align}
The odd-index congruence $\psi_A(u_0)\equiv u_0\pmod\Delta$
makes $\delta$ integral; $u_0\ge3$ and Pell growth make it
positive. Congruence~\eqref{eq:powercong} makes $\rho$ integral.
Also $E_A(u_0)>\psi_A(u_0)>2^{u_0}$, so $\rho>0$.

To prove the remaining positivity, observe that
\[
 H\sigma=E_A(R)-E_A(u_0)-X+W_0.
\]
Using $R\ge u_0+2$ and
$\psi_A(j)<E_A(j)\le2\psi_A(j)$ for the relevant positive
indices,
\[
 E_A(R)-E_A(u_0)
 >\psi_A(u_0+2)-2\psi_A(u_0)
 >(4A^2-2A-3)\psi_A(u_0)>A>X.
\]
Hence $\sigma>0$. The restored parent input quantities satisfy
\[
 \kappa=u_0+\delta\Delta,
 \qquad \mu_0=W_0+a\kappa+\rho H,
 \qquad \mu_0^2-\Delta\kappa^2=1.
\]
All its other factors are identical, and $\rho+\sigma=\gamma$.
Thus every positive zero of $U$ gives a positive zero of the exact
parent polynomial at $x=1$. We may now apply that parent's complete
ordinary-input acceptance theorem.

Conversely a positive parent zero at $x=1$ has $W=W_0$ by its
proved input-decoding theorem. Erase $\delta,\rho,\sigma$,
replace the latter two by $\gamma=\rho+\sigma$, and use the
affine guard. Every retained factor is one, giving $U=0$.
The parent's input-rank theorem makes the restoration unique.
This establishes the full positive-zero correspondence.

Finally $\Delta=(a+1)(a+3)>0$ on the entire positive supplied
domain. Hence $P=\Delta U$ has precisely the same positive zeros.
Combining the inherited compiler, Lemma~\ref{lem:normalize},
and the K215 calibration proves Theorem~\ref{thm:main}.

\section{Operation counts and exact degrees}\label{sec:counts}

\subsection{The complete evaluator}

The normalized parent evaluator uses $85=48\mathrm M+37\mathrm A$
operations. Replacing $\rho+\sigma$ by the supplied $\gamma$
removes one addition. Fixing $x=1$ removes the multiplication
$2d\,x$. Its eleven input-loader instructions are replaced by
two instructions computing $3W+c_0$. The resulting unscaled
circuit has $74=42\mathrm M+32\mathrm A$ operations.

The exact coefficient sharing from \cite{Scaled84} saves a further
multiplication. The unscaled strong/auxiliary block computes
$S=\Delta i^2c^4$ and $A_0=\Delta S$. Instead compute
\[
 v_0=i(\Delta c^2),\qquad A_0=v_0^2,
 \qquad \widetilde N_6=\Delta f^2-A_0=\Delta N_6.
\]
The quantity $\Delta c^2$ is already paid for in the main norm.
Use $\widetilde N_6$ in the final product and subtract $\Delta$
instead of one. This is exactly $P=\Delta U$ over every
commutative ring, not just on solutions. The complete circuit,
including its finalizer, therefore has $73$ gates.

For the system, remove the six product multiplications and final
subtraction from the $74$-gate unscaled evaluator, and remove the
two-gate affine guard. This leaves $65=35\mathrm M+30\mathrm A$
operations producing the seven comparison pairs. Computing the
seven residuals, their squares, and their sum adds $20$ operations,
so the direct SOS evaluator has $85=42\mathrm M+43\mathrm A$
operations. Equalities are counted separately in the comparison
model; their arithmetic cost is fully charged in the SOS version.

\subsection{Uniform exact degree}

Assign degree one to each supplied unknown and degree zero to
every fixed compiler numeral. Put
\[
 Q=\beta J,\quad k_0=\eta+\zeta,\quad
 C_1=Q-F-Z-\alpha,\quad W_1=Q-F-2Z-\alpha,
\]
and $T_1^{\rm tr}=wC_1-tQ$. The respective factor degrees of
$U+1$ are
\[
 22,\ 18,\ 1,\ 60,\ 7,\ 2,\ 34,
\]
whose sum is $144$. The main norm's apparent top terms cancel:
\[
 D^2-\Delta c^2=(X+\gamma H)^2+2ac(X+\gamma H)-Hc^2.
\]
Its degree is $18$, not the larger bound obtained by blindly
propagating gate degrees.

Multiplying the leading forms gives
\begin{equation}\label{eq:Uleader}
 U_{144}=-24Q^{83}h\gamma i^4 k_0^{13}w^{11}s^{24}
 T^2f^2 W_1 T_1^{\rm tr}.
\end{equation}
Since $\Delta_{12}=w^2s^2Q^8$, the scaled leader is
\begin{equation}\label{eq:Pleader}
 P_{156}=-24Q^{91}h\gamma i^4 k_0^{13}w^{13}s^{26}
 T^2f^2 W_1 T_1^{\rm tr}.
\end{equation}
For example the coefficient of
\[
 J^{92}h\gamma i^4\eta^{13}w^{13}s^{26}T^2f^2\alpha t
\]
is $-24\beta^{92}\ne0$. Thus $156$ is the exact degree
uniformly for every admissible fixed compiler, not a sample-based
estimate. The analogous argument gives exact degree $144$ for $U$.

The largest comparison residual degree is $60$, from the auxiliary
norm. Its square supplies the unique degree-$120$ part of the
direct SOS. These claims are independently corroborated by full
dense univariate coefficient executions over two finite fields;
those computations support, and do not replace, the leading-form
proof.

\section{Unrestricted integers and a small-degree system}\label{sec:integer}

The positive domain in the principal theorem must not be silently
changed to all integers. In particular $\Delta=0$ creates
off-domain zeros of $P=\Delta U$ which have no computational
meaning. We give explicit domain conversions.

\subsection{One equation in forty-four unrestricted integers}

The coordinates $\tau,y$ occur only squared. At a solution of
\eqref{eq:system}, $\tau\ne0$ because
$\tau^2=1+L(L+k)>1$, and $y\ne0$ because otherwise
$A_0V^2=1$ with $A_0>1$. Consequently these two roots may range
over all integers. The unit-sign proof remains valid for signed
roots: use $|\tau|$ in the negative-Pell argument and retain the
same square congruences. Thus every zero on this partially signed
domain has nonzero roots, and absolute values restore a positive solution.
The remaining fourteen coordinates must stay positive.

Every positive integer has a representation
\begin{equation}\label{eq:positive3}
 v=1+a^2+b^2+c(c+1),\qquad a,b,c\in\Z.
\end{equation}
Indeed $4(v-1)+1$ is not of the excluded three-square form
$4^r(8s+7)$. In a three-square representation it has exactly
one odd square modulo four, giving \eqref{eq:positive3}.
Conversely its right-hand side is always positive for integer
$a,b,c$. Substitute \eqref{eq:positive3} into the fourteen
positive coordinates and keep $\tau,y$ signed. This defines a
single polynomial in $14\cdot3+2=44$ unrestricted integer
unknowns, with exactly the required solvability equivalence.
Its degree is $312$. Restricting all the $b,c$ parameters to
zero verifies noncancellation of the doubled leading degree.
The literal substitution costs $14(3\mathrm M+4\mathrm A)$,
giving $171=83\mathrm M+88\mathrm A$ total operations.

\subsection{Forty-six quadratic equations in eighty-three integers}

Starting from the $65$-operation comparison circuit, inline every
affine gate. Three multiplication gates are affine as well:
$\beta J$, $4a$, and $q\widehat M_F$. There remain $32$
genuinely nonlinear multiplication-register equations.
Every remaining comparison is affine in the supplied coordinates
and these product registers. Solve the seven comparisons using
the pivots
\[
 \tau^2\text{-register},\ D^2\text{-register},\ f^2\text{-register},
 \text{auxiliary-product register},\ hE\text{-register},
 t(q-1)\text{-register},\ \alpha.
\]
Each pivot has coefficient $1$ or $-1$. Substitution therefore
preserves integer solutions and introduces no divisibility condition.
All product equations remain quadratic.

For each of the fourteen positive coordinates use the particularly
simple positivity equation
\begin{equation}\label{eq:positiveeq}
 4v-3=a_v^2+b_v^2+c_v^2.
\end{equation}
It is solvable exactly when $v\ge1$, by the same three-square
theorem; $\alpha$ is replaced by its already solved affine
expression. The two signed roots need no positivity equation.
The total is $32+14=46$ quadratic equations. There are
$(16+32-7)+42=83$ unrestricted integer variables.

The archive lists all $46$ residual polynomials explicitly in
\texttt{data/quadratic83.txt}, and serializes every coefficient
and exponent in \texttt{data/quadratic83.json}. Appendix~\ref{app:quadratic}
prints their defining equations. Let them be $Q_1,\ldots,Q_{46}$.
Then the single equation
\begin{equation}\label{eq:quartic}
 \boxed{\qquad Q_1^2+\cdots+Q_{46}^2=0\qquad}
\end{equation}
has exactly the same integer solvability condition. It has exact
degree four: the fourth power of the first positivity parameter
for $J$ has coefficient one. This quartic is a useful alternative
when low degree matters more than a short evaluator or few variables.

\section{Why the equivalence is a theorem of ZFC}\label{sec:internal}

The proof chain is
\[
 \begin{aligned}
 \exists\mathbf x\;P(\mathbf x)=0
 &\Longleftrightarrow\text{the fixed compiled input is accepted}\\
 &\Longleftrightarrow\operatorname{Halt}(\mathrm{K215},\mathrm{blank})\\
 &\Longleftrightarrow\exists n\;\operatorname{Proof}_{\ZFC}(n,\bot).
 \end{aligned}
\]
Each step has both directions. The coefficient recipe consists of
finite sets, bounded enumerations, integer sums and products, and
primitive finite-table constructions. The normalization and tableau
arguments are induction on finite computations. The Pell lemmas,
integer-unit arguments, and three-square theorem are ordinary
number theory available in ZFC. The inherited K215 calibration has
an ordinary ZFC-formalizable mathematical proof, separately supported
by its external proof-assistant endpoint \cite{K215,Endpoint}.
Thus these are proofs ZFC can carry out; their composition does not
assume the consistency or soundness of ZFC.

The delivered proof is not a string of ZFC proof-code instructions,
and the new $73$-gate composition has not been formalized in Lean
or Rocq. An external Lean theorem is not, merely by being printed,
a derivation encoded in ZFC's proof syntax. The present claim uses
the ordinary mathematical arguments and their formalizability,
and explicitly records which substantial inherited lemmas it uses.

Negating the existential equivalence gives the requested result.
If ZFC is consistent, the equation has no standard positive solution,
and ZFC cannot prove that lack of solutions, by its second
incompleteness theorem. We do not infer full two-sided
unprovability from consistency alone: ruling out a ZFC proof of
the false existential assertion uses the relevant soundness
assumption, such as $\Sigma_1$-soundness. The requested
nonexistence--consistency equivalence does not require such an
additional assumption.

\section{Verification, reproducibility, and remaining work}\label{sec:verify}

The package has three distinct verification layers.

\paragraph{Literal machine and compiler data.}
The normalizer authenticates the exact K215 table, checks all
$430$ inherited transitions and the complete $660$-row normalized
table, checks several input-erasure lengths, and compares $120{,}000$
paired simulation steps. These prefix checks test the implementation;
they are not evidence of unconditional nonhalting. The compiler
includes a two-sided $126$-cell accepting tableau fixture, exact
marker separation, both overlap directions, and five actually
materialized small compiler instances. A separate source audit
compares $25$ sparse layout attributes against the two pinned
ancestor compiler functions, including both optional-correction
branches. None of these tests materializes a K215 accepting run.
The factored window count was independently compared with all
$209{,}952$ weighted cross assignments for an $18$-tile fixture;
the scalar formulas agree with all five materialized compiler
fixtures. The general counting derivation is given in the coefficient note.

\paragraph{Finite algebra.}
The supplied checker verifies complete source closure, liveness,
interfaces, and operation ledgers. Eighty signed/positive arithmetic
fixtures check the all-value identity $P=\Delta U$ and the exact
SOS finalizer. Two full univariate coefficient evaluations check
the degrees $144,156,120$. The quadratic compiler performs exact
symbolic substitutions for all $32$ product equations and all
seven comparisons, verifies every variable is live, and verifies
the coefficient-one quartic witness. These are exact finite
checks, not a finite search used as a substitute for the
unbounded mathematical equivalence.

\paragraph{Mathematics.}
A separate review within this research collaboration checked unit signs, pretyping bounds,
auxiliary positivity, the exact strong-rank and step-down
hypotheses, the noncircular recovery $X=2^R$, and all three
restored input witnesses. The final new theorem still relies on
the explicitly cited inherited K215 and native compiler results.
Fresh full builds of their Lean and Rocq dependency trees were
not performed for this report.

From the archive directory, the standard-library checks run as follows:
\begin{lstlisting}
python3 src/verify_all.py
\end{lstlisting}
The individual entry points are also available:
\begin{lstlisting}
python3 src/diophantine.py --emit --verify
python3 src/normalize_k215_direct.py data/k215.tm --verify
python3 src/k215_constants.py data/normalized_k215.json --self-test
python3 src/quadratic_elimination.py --expect data/quadratic83.json
python3 src/audit_compiler_transcription.py
\end{lstlisting}
The TeX source compiles with two passes of pdfLaTeX. The archive
also contains the generated JSON circuits, quadratic equations,
verification receipts, source provenance, and a build script.
The ordinary commands above are bounded practical checks; the
opt-in full K215 coefficient expansion is excluded from them.

The next useful formalization step is to state a generic
fixed-marker theorem for the native parent and prove the affine
unit-guard reduction once. This avoids recompiling or re-proving
the entire machine simulation for each arithmetic simplification.
Coefficient height is a separate research objective: a compact
direct register-machine or proof-calculus compiler could produce
much smaller constants without improving the $73$-gate count.
The present package makes that boundary visible rather than
hiding it in an unspecified universal index.

\appendix
\section{The explicit quadratic equations}\label{app:quadratic}

All symbols in the following listing are unrestricted integer
unknowns except the six fixed coefficients
\texttt{Bm1}, \texttt{Kconstant}, \texttt{twice\_cell\_bits},
\texttt{MC}, \texttt{MF}, and \texttt{W0}. They mean respectively
$\beta,K_0,2d,M_C,\widehat M_F,W_0$ from
Section~\ref{sec:constants}. Each expression is set equal to zero.
The labels at the left identify eliminated source-register
equations; a label is not an additional variable. A machine-readable
complete variable list accompanies the equations in the archive.
\lstinputlisting[firstline=10]{data/quadratic83.txt}

\begin{thebibliography}{99}
\raggedright
\bibitem{K215} ProveIt, \emph{Formal correctness of the exact
215-state machine}, snapshot \snap.
\href{https://github.com/VladimirReshetnikov/ProveIt/blob/cdd582bec13f908df3711fa9e05d1ddac368a5db/Computability/BusyBeaver/Research/ZFC/k225-formal/verification/k215/README.md}{Pinned K215 verification overview}.
\bibitem{Endpoint} ProveIt, \emph{KEndpoint215.lean}, same snapshot.
\href{https://github.com/VladimirReshetnikov/ProveIt/blob/cdd582bec13f908df3711fa9e05d1ddac368a5db/Computability/BusyBeaver/Lean/BusyBeaver/ZFC/KEndpoint215.lean}{Literal Lean endpoint and dependency composition}.
\bibitem{Raw75} ProveIt, \emph{A 75-operation fixed-index universal
certificate}, \texttt{FIXED\_RAW\_UNIVERSAL\_75\_PROOF.md}, same snapshot.
\href{https://github.com/VladimirReshetnikov/ProveIt/blob/cdd582bec13f908df3711fa9e05d1ddac368a5db/Computability/HilbertTenthProblem/Papers/1980/FIXED_RAW_UNIVERSAL_75_PROOF.md}{Full native mathematical proof}.
\bibitem{Compiler75} ProveIt, \emph{Modified half-binomial compiler},
\texttt{complete75\_half\_binomial\_compiler.md} and its Python source,
same snapshot.
\href{https://github.com/VladimirReshetnikov/ProveIt/blob/cdd582bec13f908df3711fa9e05d1ddac368a5db/Computability/HilbertTenthProblem/Papers/research-wip/native-stream-queue/complete75_half_binomial_compiler.md}{Compiler proof and fixed-numeral conventions}.
\bibitem{Bezout85} ProveIt, \emph{An auxiliary quotient gives a complete
universal polynomial in 85 operations},
\texttt{complete85\_auxiliary\_bezout\_projection.md}, same snapshot.
\href{https://github.com/VladimirReshetnikov/ProveIt/blob/cdd582bec13f908df3711fa9e05d1ddac368a5db/Computability/HilbertTenthProblem/Papers/research-wip/native-stream-queue/complete85_auxiliary_bezout_projection.md}{Normalized parent, positivity, and input interface}.
\bibitem{Scaled84} ProveIt, \emph{A complete universal polynomial in
84 operations}, \texttt{complete84\_scaled\_strong\_output.md}, same snapshot.
\href{https://github.com/VladimirReshetnikov/ProveIt/blob/cdd582bec13f908df3711fa9e05d1ddac368a5db/Computability/HilbertTenthProblem/Papers/research-wip/native-stream-queue/complete84_scaled_strong_output.md}{Exact positive scaling and coefficient sharing}.
\bibitem{Relaxed} ProveIt, \emph{Relaxed auxiliary Pell proof},
\texttt{PELL\_RELAXED\_AUXILIARY\_PROOF.md}, same snapshot.
\href{https://github.com/VladimirReshetnikov/ProveIt/blob/cdd582bec13f908df3711fa9e05d1ddac368a5db/Computability/HilbertTenthProblem/Papers/1980/PELL_RELAXED_AUXILIARY_PROOF.md}{Local strong-rank lemma}.
\bibitem{Reverse89} ProveIt, \emph{Reversed auxiliary refinement},
\texttt{complete75\_reversed\_auxiliary89.md}, same snapshot.
\href{https://github.com/VladimirReshetnikov/ProveIt/blob/cdd582bec13f908df3711fa9e05d1ddac368a5db/Computability/HilbertTenthProblem/Papers/research-wip/native-stream-queue/complete75_reversed_auxiliary89.md}{Signed auxiliary step-down argument}.
\bibitem{Jones82} J. P. Jones, \emph{Universal Diophantine equation},
Journal of Symbolic Logic \textbf{47} (1982), 549--571.
\href{https://doi.org/10.2307/2273588}{doi:10.2307/2273588}.
\bibitem{BayerDavid} J. Bayer and M. David,
\emph{A Formal Proof of Complexity Bounds on Diophantine Equations},
ITP 2025, LIPIcs 352, Article 3.
\href{https://doi.org/10.4230/LIPIcs.ITP.2025.3}{doi:10.4230/LIPIcs.ITP.2025.3}.
\end{thebibliography}
\end{document}

